\documentclass[../Main.tex]{subfiles}
\begin{document}
\section{Research Goal, Questions, and Methods}
In this project, the Goal, Question, Metrics (GQM) approach \citep{basili_software_1992} is used to define the software metrics. This methodology introduces a conceptual, operational, and quantitative level. At the conceptual level, a goal for what the project should achieve is defined. The operational level includes a set of research questions associated with the goal, which can indicate whether or not the goal is accomplished. The quantitative level defines a set of metrics that will help answer the research questions.

\subsection{Research Goal}
\label{sec:research_goal}
The research goal of this project is:
\textit{The goal of this project is to develop one or more game concepts for an alternative game that has a socially beneficial effect, and that can be developed, explored, and tested.}

\subsection{Research Questions}
\label{sec:research_questions}
Four research questions are introduced to concretize the research goal:

\textbf{RQ1}: \textit{What makes a game engaging?} \\
This research question will investigate theories and concepts used in game development that makes games engaging. The research will explore concepts that can enhance the players’ enjoyment, motivation, and engagement.

\textbf{RQ2}: \textit{What is an alternative game?} \\
This research question investigates what alternative games are and how they differ from other games.

\textbf{RQ3}: \textit{How can alternative games be categorized and conceptualized?} \\
This research question will examine how existing alternative games can be categorized using a structured approach. The research will highlight existing alternative games and point to what exactly makes them alternative. It will also investigate if a structured approach improves the process of developing new alternative game concepts.

\textbf{RQ4}: \textit{Which alternative game concepts can be developed with existing technology that has a socially beneficial purpose?} \\
This research question will consider game concepts that can be developed based on the knowledge gained from literature and technology reviews. The concepts that we will develop should comply with the definition of alternative games and have a socially beneficial purpose. It is also of interest to study how these concepts can contribute to the study of alternative games through our interpretation of alternative games. The most promising concepts will be conceptualized extensively and developed (at a later stage). 

\subsection{Research Methods}
The research questions above will be answered through an analysis of relevant literature. These areas include game development, game flow, existing alternative games, game platforms, and available game technology. \citet{aveyard_doing_2014} states that a literature review is a “comprehensive study and interpretation of literature that relates to a particular topic.” The insight from this literature study on existing games, relevant technologies, and game development concepts will be summarized in the prestudy. This part should provide the background information needed to do further research in this field. Through the review of the literature, we will seek to answer the research questions. The answers will give insight into what alternative games are, which alternative games already exist, and how they contribute to the field. Also, relevant game design theories on how to create fun and engaging games will be reviewed. 

The research field of alternative games covers many different types of games (the pool of existing examples are very diverse). Insight acquired from studying what makes games engaging in \textbf{RQ1} and existing games in \textbf{RQ3} will be used to develop a model for categorizing alternative games. The structured modeling approach will be used to find synergies between the games and tested during the development of new concepts, as described in \textbf{RQ4}. The most promising concepts will be evaluated further before one final concept is prototyped.


\biblio % Needed for referencing to working when compiling individual subfiles - Do not remove
\end{document}